% Part: normal-modal-logic
% Chapter: completeness
% Section: frame-completeness

\documentclass[../../../include/open-logic-section]{subfiles}

\begin{document}

\olfileid{nml}{com}{fra}

\olsection{框架完备性}

只要证明给定逻辑的典范模型具有相应的框架性质，\Log{K} 的完备性定理便可推广到其他模态系统。

\begin{thm}\ollabel{thm:completeframeprops}
  若正规模态逻辑 $\Sigma$ 包含 \olref{tab:correspondencetable} 左侧列出的公式之一，
  则 $\Sigma$ 的典范模型具有右侧对应的性质。
\end{thm}

\begin{table}[htp]
  \centering
    \begin{tabular}{| l || l |}
      \hline
      \emph{若 $\Sigma$ 包含……} & \emph{……则 $\Sigma$ 的典范模型是：} \\
      \hline \hline
      \Ax{D}:\;\; $\Box!A \lif \Diamond !A$ & \quad 持续的； \\
      \hline
      \Ax{T}:\;\; $\Box!A \lif !A$ &\quad  自反的；\\
      \hline
      \Ax{B}: \;\;$!A \lif \Box\Diamond!A$ &\quad  对称的； \\
      \hline
      \Ax{4}: \;\; $\Box!A \lif \Box\Box!A$ & \quad  传递的； \\
      \hline 
      \Ax{5}: \;\; $\Diamond !A \lif \Box\Diamond!A$& \quad  欧几里得的。\\
      \hline
    \end{tabular}
    \caption{基本对应事实。}
    \ollabel{tab:correspondencetable}
\end{table}

\begin{proof}
  我们依次处理每一种情况。

  假设 $\Sigma$ 包含 \Ax{D}，并设 $\Delta \in W^\Sigma$；需证存在 $\Delta'$ 使得 $R^\Sigma
  \Delta\Delta'$。只需证 $\Box^{-1}\Delta$ 是 $\Sigma$-一致的，因为由此根据 Lindenbaum 引理，存在完全的 $\Sigma$-一致集 $\Delta' \supseteq
  \Box^{-1}\Delta$，而根据 $R^\Sigma$ 的定义\iftag{prvBox}{}{及 \olref[mod]{lem:box-iff-diamond} }可得
  $R^\Sigma \Delta\Delta'$。反设 $\Box^{-1}\Delta$ 是 $\Sigma$-\emph{不一致}的，即
  $\Box^{-1}\Delta \Proves[\Sigma] \lfalse$。根据 \olref[mod]{lem:box2}，
  $\Delta \Proves[\Sigma] \Box\lfalse$；又因 $\Sigma$ 包含
  \Ax{D}，亦有 $\Delta \Proves[\Sigma] \Diamond\lfalse$。但 $\Sigma$
  是正规的，故 $\Sigma \Proves \lnot\Diamond \lfalse$
  (\olref[prf][nor]{prop:notDiamondBot})，从而 $\Delta
  \Proves[\Sigma] \lnot\Diamond \lfalse$，与 $\Delta$ 的一致性矛盾。

  假设 $\Sigma$ 包含 \Ax{T}，并设 $\Delta \in
  W^\Sigma$。欲证 $R^\Sigma \Delta\Delta$，
  即\iftag{prvBox}{}{根据 \olref[mod]{lem:box-iff-diamond}，}
  $\Box^{-1}\Delta \subseteq \Delta$。但若 $\Box!A \in \Delta$，则由 \Ax{T} 亦有 $!A \in \Delta$，如所求。

  假设 $\Sigma$ 包含 \Ax{B}，并设 $R^\Sigma
  \Delta\Delta'$，其中 $\Delta$, $\Delta' \in W^\Sigma$。需证
  $R^\Sigma \Delta'\Delta$，即\iftag{prvBox}{
    $\Box^{-1}\Delta' \subseteq \Delta$。根据
    \olref[mod]{lem:box-iff-diamond}，这等价于}{}
  $\Diamond\Delta \subseteq \Delta'$。故设 $!A \in \Delta$。由
  \Ax{B}，亦有 $\Box\Diamond!A \in \Delta$。由 $R^\Sigma \Delta\Delta'$ 的假设\iftag{prvBox}{}{及
    \olref[mod]{lem:box-iff-diamond}}，可得 $\Box^{-1}\Delta
  \subseteq \Delta'$，从而 $\Diamond!A \in \Delta'$，如所求。

  假设 $\Sigma$ 包含 \Ax{4}，并设 $R^\Sigma
  \Delta_1\Delta_2$ 且 $R^\Sigma \Delta_2\Delta_3$。需证
  $R^\Sigma \Delta_1\Delta_3$。由假设\iftag{prvBox}{}{及
    \olref[mod]{lem:box-iff-diamond}}可得 $\Box^{-1}\Delta_1
  \subseteq \Delta_2$ 与 $\Box^{-1}\Delta_2 \subseteq \Delta_3$。为证 $R^\Sigma \Delta_1\Delta_3$，只需证
  $\Box^{-1}\Delta_1 \subseteq \Delta_3$。故设 $!B \in
  \Box^{-1}\Delta_1$，即 $\Box!B \in \Delta_1$。由 \Ax{4}，亦有
  $\Box\Box!B \in \Delta_1$，由假设首先可得 $\Box!B \in \Delta_2$，进而得 $!B \in \Delta_3$，如所求。

  假设 $\Sigma$ 包含 \Ax{5}，并设 $R^\Sigma
  \Delta_1\Delta_2$ 且 $R^\Sigma \Delta_1\Delta_3$。需证
  $R^\Sigma \Delta_2\Delta_3$。由第一个假设可得 $\Box^{-1}\Delta_1 \subseteq \Delta_2$\iftag{prvBox}{}{（根据
    \olref[mod]{lem:box-iff-diamond}）}，且第二个假设等价于 $\Diamond\Delta_3 \subseteq \Delta_1$\iftag{prvBox}{（根据 \olref[mod]{lem:box-iff-diamond}）}{}。要证 $R^\Sigma
  \Delta_2\Delta_3$\iftag{prvBox}{，根据
    \olref[mod]{lem:box-iff-diamond}，只需}{，需}证 $\Diamond\Delta_3 \subseteq \Delta_2$。设 $\Diamond!A \in
  \Diamond\Delta_3$，即 $!A \in \Delta_3$。由第二个假设，$\Diamond!A \in \Delta_1$，再由 \Ax{5}，$\Box\Diamond!A \in
  \Delta_1$。而由第一个假设可得 $\Diamond!A \in \Delta_2$，如所求。
  \end{proof}

作为推论，我们得出若干系统的完备性结果。例如，可知 $\Log{S5} = \Log{KT5} =
\Log{KTB4}$ 相对于所有自反且欧几里得的模型类是完备的，而该类即所有自反、对称且传递的模型类。

\begin{thm}\ollabel{thm:generaldet}
  令 $\mClass{C}_\Ax{D}$、$\mClass{C}_\Ax{T}$、
  $\mClass{C}_\Ax{B}$、$\mClass{C}_\Ax{4}$ 和
  $\mClass{C}_\Ax{5}$ 分别为所有持续的、自反的、
  对称的、传递的和欧几里得的模型类。那么，对于
  公理 \Ax{D}、\Ax{T}、\Ax{B}、\Ax{4} 和 \Ax{5} 中的任意
  模式 $!A_1$, \dots, $!A_n$，系统
  $\Log{K}!A_1 \dots !A_n$ 都由模型类
  $\mClass{C} = \mClass{C}_{!A_1} \cap \dots
  \cap \mClass{C}_{!A_n}$ 所决定。
\end{thm}

\begin{prop}
  设 $\Sigma$ 为正规模态逻辑，则：
  \begin{enumerate}
  \item\ollabel{prop:anotherfive-a}%
    若 $\Sigma$ 包含模式 $\Diamond!A \lif \Box 
    !A$，则 $\Sigma$ 的典范模型是部分函数的。 
  \item 若 $\Sigma$ 包含模式 $\Diamond!A \liff \Box 
    !A$，则 $\Sigma$ 的典范模型是函数的。 
  \item 若 $\Sigma$ 包含模式 $\Box\Box!A \lif \Box 
    !A$，则 $\Sigma$ 的典范模型是弱稠密的。 
  \end{enumerate}
（这些框架性质的定义参见 \olref[frd][acc]{tab:anotherfive}）。
\end{prop}

\begin{proof}
  \begin{enumerate}
  \item 假设 $\Sigma$ 包含模式 $\Diamond !A \lif
    \Box !A$，欲证 $R^\Sigma$ 是部分函数的，需证对于任意 $\Delta_1$, $\Delta_2$, $\Delta_3 \in
    W^\Sigma$，若 $R^\Sigma \Delta_1\Delta_2$ 且 $R^\Sigma
    \Delta_1\Delta_3$，则 $\Delta_2=\Delta_3$。由 $R^\Sigma
    \Delta_1\Delta_2$ 可得 $\Box^{-1}\Delta_1 \subseteq \Delta_2$，
    由 $R^\Sigma \Delta_1\Delta_3$ 亦得 $\Box^{-1}\Delta_1
    \subseteq \Delta_3$\iftag{prvBox}{}{（两者均依据
      \olref[mod]{lem:box-iff-diamond}）}。若能证得 $\Delta_2 \subseteq \Delta_3$ 与 $\Delta_3 \subseteq
    \Delta_2$ 两个包含关系，则等式 $\Delta_2=\Delta_3$ 成立。为证第一个包含关系，设 $!A \in \Delta_2$；则
    $\Diamond!A \in \Delta_1$，由该模式及 $\Delta_1$ 的演绎封闭性，可得 $\Box!A \in \Delta_1$，再由假设 $R^\Sigma \Delta_1\Delta_3$ 可得 $!A \in \Delta_3$。第二个包含关系的证明类似。

  \item 这直接由 \olref{prop:anotherfive-a} 及 \olref{thm:completeframeprops} 中的持续性证明得出。

  \item 假设 $\Sigma$ 包含模式 $\Box\Box!A \lif \Box!A$，欲证 $R^\Sigma$ 是弱稠密的。设 $R^\Sigma
    \Delta_1\Delta_2$。需证存在完全的 $\Sigma$-一致集 $\Delta_3$，使得 $R^\Sigma
    \Delta_1\Delta_3$ 且 $R^\Sigma \Delta_3\Delta_2$。令：
    \[
    \Gamma = \Box^{-1}\Delta_1 \cup \Diamond\Delta_2.
    \]
    只需证 $\Gamma$ 是 $\Sigma$-一致的，因为由此根据 Lindenbaum 引理，它可扩充为完全的
    $\Sigma$-一致集~$\Delta_3$，满足 $\Box^{-1}\Delta_1
    \subseteq \Delta_3$ 且 $\Diamond\Delta_2 \subseteq \Delta_3$，
    即 $R^\Sigma \Delta_1\Delta_3$\iftag{prvBox}{}{（依据
      \olref[mod]{lem:box-iff-diamond}）}且 $R^\Sigma
    \Delta_3\Delta_2$\iftag{prvBox}{（依据
      \olref[mod]{lem:box-iff-diamond}）}{}。

    反设 $\Gamma$ 不是一致的。则
    存在公式 $\Box!A_1$, \dots, $\Box!A_n \in \Delta_1$
    及 $!B_1$, \dots,~$!B_m \in \Delta_2$，使得 \[!A_1, \dots,
    !A_n, \Diamond!B_1, \dots, \Diamond!B_m \Proves[\Sigma]
    \lfalse.\] 由于 $\Diamond (!B_1 \land \dots \land !B_m) \to
    (\Diamond!B_1 \land \dots \land \Diamond!B_m)$ 在
    任何正规模态逻辑中皆可推导，我们作如下推演，从而与 $\Delta_2$ 的一致性矛盾：
    \begin{align*}
      !A_1, \dots, !A_n, & \Diamond!B_1, \dots,
      \Diamond!B_m \Proves[\Sigma] \lfalse \\
      !A_1, \dots,!A_n & \Proves[\Sigma] (\Diamond!B_1 \land
      \dots \land \Diamond!B_m) \lif \lfalse\\
      & \qquad\text{由演绎定理}\\
      & \qquad\text{\olref[prf][prp]{prop:derivabilityfacts}\olref[prf][prp]{prop:derivabilityfacts-deduction}，以及 \Taut} \\
      !A_1, \dots,!A_n & \Proves[\Sigma]
      \Diamond(!B_1 \land
      \dots \land !B_m) \lif \lfalse\\
      & \qquad \text{因为 $\Sigma$ 是正规的} \\
      !A_1, \dots,!A_n & \Proves[\Sigma]
      \lnot\Diamond (!B_1 \land \dots \land !B_m)\\
      & \qquad\text{由 \PL} \\
      !A_1, \dots,!A_n & \Proves[\Sigma]
      \Box\lnot (!B_1 \land \dots \land !B_m)\\
      & \qquad\text{对 $\lnot\Diamond$ 使用 $\Box\lnot$} \\
      \Box!A_1, \dots,\Box!A_n
      & \Proves[\Sigma] \Box\Box \lnot (!B_1 \land \dots \land !B_m)\\
      & \qquad\text{由 \olref[mod]{lem:box1}} \\
      \Box!A_1, \dots,\Box!A_n
      & \Proves[\Sigma]
      \Box\lnot (!B_1 \land \dots \land !B_m)\\
      &\qquad\text{由模式 $\Box\Box!A \lif \Box!A$} \\
      \Delta_1 & \Proves[\Sigma]
      \Box\lnot (!B_1 \land \dots \land !B_m)\\
      &\qquad\text{由单调性，\olref[prf][prp]{prop:derivabilityfacts}\olref[prf][prp]{prop:derivabilityfacts-monotonicity}} \\
      & \Box\lnot (!B_1 \land \dots \land !B_m) \in \Delta_1\\
      &\qquad\text{由演绎封闭性}; \\
      & \lnot (!B_1 \land \dots \land !B_m) \in \Delta_2\\
      &\qquad \text{因为 }
      R^\Sigma \Delta_1\Delta_2.
    \end{align*}
  \end{enumerate}
\end{proof}

基于这些例子，人们或许会认为每个模态逻辑系统~$\Sigma$ 都是\emph{完备的}，即：若一个公式在每个使得 $\Sigma$ 的所有定理皆有效的框架中都有效，则 $\Sigma$ 能证明该公式。遗憾的是，有许多系统在这种意义上并不完备。

\end{document}